\documentclass[10pt]{article}
\usepackage[a4paper,margin=23mm]{geometry}
\usepackage{amsmath,amssymb}
\usepackage{hyperref}
\usepackage{url}
\setlength{\emergencystretch}{3em}
\title{A critical equality graph refutes Conjecture 00000003483}
\author{gaochengzhecpu (AI-assisted submission)}
\date{3 October 2026}
\begin{document}
\maketitle

\paragraph{Statement and definitions.}
The conjecture asserts that equality in the edge-count expression
$e=(5v-2)/3$ for four-critical graphs is attained only by nearly-four-regular
double-critical graphs. We refute the necessary double-critical part of this assertion.
A graph is four-critical if it has chromatic number $4$ and every proper subgraph is $3$-colorable.
A connected $k$-chromatic graph is double-critical if deleting both endpoints of
every edge lowers the chromatic number by two, equivalently leaves a $(k-2)$-colorable
graph; see Kawarabayashi, Pedersen, and Toft~\cite{KPT}.
The equivalence follows because restoring the two deleted vertices uses at most two new colors.

\paragraph{The graph and equality.}
Let $G$ have vertex set $\{0,1,\ldots,6\}$ and the following eleven unordered edges:
\[
02,\ 03,\ 05,\ 06,\ 12,\ 13,\ 14,\ 23,\ 45,\ 46,\ 56.
\]
This is the Haj\'os join of two copies of $K_4$: delete edges $01$ and $04$
from the two four-vertex graphs, identify their vertices labelled $0$,
and add the edge $14$.
There are $v=7$ vertices and $e=11$ edges, so
\[
3e=33=5v-2.
\]

\paragraph{Four-criticality.}
A proper four-coloring, in vertex order $0,\ldots,6$, is
$(0,0,1,2,1,2,3)$.
There is no proper three-coloring: vertices $2,3$ are adjacent and both are
adjacent to $0,1$, forcing $0$ and $1$ to have the same third color.
Similarly, vertices $5,6$ are adjacent and both adjacent to $0,4$,
forcing $0$ and $4$ to have the same color.
This contradicts the edge $14$.
Thus $\chi(G)=4$.

For a fully explicit criticality check, the following tables give proper
three-colorings after every single vertex or edge deletion.
Entries always list colors at vertices $0,\ldots,6$; the color at a deleted
vertex is ignored.
\[
\begin{array}{c|c@{\qquad}c|c}
\text{deleted vertex}&\text{colors}&\text{deleted edge}&\text{colors}\\ \hline
0&(0,0,1,2,1,0,2)&02&(0,1,0,2,0,1,2)\\
1&(0,0,1,2,0,1,2)&03&(0,1,2,0,0,1,2)\\
2&(0,1,0,2,0,1,2)&05&(0,0,1,2,1,0,2)\\
3&(0,1,2,0,0,1,2)&06&(0,0,1,2,1,2,0)\\
4&(0,0,1,2,0,1,2)&12&(0,1,1,2,0,1,2)\\
5&(0,0,1,2,1,0,2)&13&(0,1,2,1,0,1,2)\\
6&(0,0,1,2,1,2,0)&14&(0,0,1,2,0,1,2)\\
&&23&(0,1,2,2,0,1,2)\\
&&45&(0,0,1,2,1,1,2)\\
&&46&(0,0,1,2,1,2,1)\\
&&56&(0,0,1,2,1,2,2)
\end{array}
\]
Every proper subgraph omits a vertex or omits an edge; restrict the
corresponding listed coloring. Thus every proper subgraph is $3$-colorable
and $G$ is four-critical in the full, standard sense.

\paragraph{Failure of double-criticality.}
The vertices $0$ and $2$ are adjacent. After deleting them, the vertices
$4,5,6$ still induce a triangle. Hence
$\chi(G-\{0,2\})\ge3$.
Restricting the displayed coloring for $G-0$ also gives a three-coloring
of $G-\{0,2\}$, so in fact
\[
 \chi(G-\{0,2\})=3\ne4-2.
\]
Therefore $G$ is not double-critical, despite being four-critical and
attaining the stated equality. Whatever additional condition is intended
by nearly-four-regular, its conjunction with double-criticality cannot
hold for this graph.

\paragraph{Formal verification and semantic bridge.}
The self-contained Lean 4.19.0 file imports only \texttt{Std}.
Vertices are \texttt{Fin 7}. The graph is the displayed edge list.
\texttt{NormalizedSimple} checks that every edge is within the vertex set,
its endpoints are in increasing order, and the list has no duplicates;
hence the verified list length is the ordinary undirected edge count.

The theorem \path{G_four_critical} proves the standard chromatic-number
minimum and three-colorability of \emph{every} proper subgraph.
All colorings by three colors are covered by \texttt{tuple\_eta},
which identifies an arbitrary function on \texttt{Fin 7} with its seven
values. The vertex and edge certificates are exactly the displayed tables;
\path{proper_subgraph_bridge} proves the restriction argument for arbitrary
vertex subsets and arbitrary symmetric subgraph relations.
Lower color counts are excluded by embedding their colors into the three-color set.

\texttt{DoubleCriticalFour} includes connectivity, chromatic number four,
and the exact two-color minimum after deleting the endpoints of any edge.
\texttt{ProperPairDeleted} applies the ordinary proper-coloring condition
only to the surviving vertices; unused labels at deleted vertices do not
change colorability. The triangle argument is formalized without
enumerating graph families, and \path{G_pair_chromatic_three} verifies
the exact deletion chromatic number.
The theorem \path{conjecture3483_counterexample} negates the necessary
seven-vertex specialization of the equality characterization.
This specialization alone refutes the source's unrestricted characterization.

The proof addresses the equality assertion itself. It does not rely on the
source's background description of the edge inequality as an upper bound.
The principal axiom audits use only standard \texttt{propext} and
\texttt{Quot.sound}; there are no added axioms, admissions, or native
decision procedures.

\begin{thebibliography}{1}
\bibitem{KPT}
K.-i.~Kawarabayashi, A.~S.~Pedersen, and B.~Toft,
\emph{Double-critical graphs and complete minors},
Electron.\ J.\ Combin.\ \textbf{17} (2010), R87.
\href{https://www.combinatorics.org/ojs/index.php/eljc/article/view/v17i1r87}{doi:10.37236/359}.
\end{thebibliography}
\end{document}
